\section{PESGM-2027 P1 IEEE-123 PV Augmentation}
\label{sec:p1_ieee123_pv_augmentation}

\subsection{Scope and provenance}
The P1 experiment uses the canonical OpenDSS IEEE 123-bus feeder as the physical
baseline. The baseline regulator and capacitor equipment is preserved unchanged.
The only P1-specific physical augmentation is the addition of ten photovoltaic
(PV) inverter systems. The modified feeder is maintained as a separate working
copy, while the source baseline remains immutable.

The experiment instantiates the generalized controller cardinalities as
\[
N_R = 7,\qquad N_C = 4,\qquad N_P = 10.
\]

\subsection{PV composition}
The ten added PV systems are divided into
\[
N_P = N_D + N_H = 3 + 7 = 10,
\]
where $N_D=3$ are three-phase distribution-connected PV systems and $N_H=7$ are
single-phase residential or house-like PV systems.

The distribution PV buses are
\[
\mathcal B_D = \{47,76,105\},
\]
and the residential PV sites are
\[
\mathcal B_H =
\{20{\rm A},38{\rm B},41{\rm C},56{\rm B},70{\rm A},103{\rm C},113{\rm A}\}.
\]
The resulting residential phase count is three phase-A, two phase-B, and two
phase-C systems.

The placement deliberately avoids co-location with canonical regulator and
capacitor locations. The intent is to avoid a trivial local pairing between each
new DER and an existing voltage-control device, while distributing DER actuation
across electrically distinct feeder regions.

\subsection{PV penetration and nameplate sizing}
The canonical IEEE-123 peak active demand used for experiment sizing is
\[
P_L^{\rm peak}=3490~{\rm kW}.
\]
Each distribution PV is assigned
\[
P_D^{\rm rated}=210~{\rm kW},\qquad
S_D^{\rm rated}=231~{\rm kVA},
\]
and each residential PV is assigned
\[
P_H^{\rm rated}=10~{\rm kW},\qquad
S_H^{\rm rated}=11~{\rm kVA}.
\]
Hence
\[
P_{\rm PV,tot}=3(210)+7(10)=700~{\rm kW},
\]
and the active-power penetration is
\[
\gamma_{\rm PV}
=
\frac{P_{\rm PV,tot}}{P_L^{\rm peak}}
=
\frac{700}{3490}
\approx 0.200573
\approx 20.06\%.
\]
A uniform inverter oversizing ratio is used:
\[
S^{\rm rated}=1.10P^{\rm rated}.
\]

\subsection{Reactive-power capability}
PV active power is exogenous to P1 control:
\[
P_{p,t}^{\rm PV}=P_{p,t}^{\rm PV,av}.
\]
The controller chooses only reactive power, subject to
\[
\left(P_{p,t}^{\rm PV,av}\right)^2+
\left(q_{p,t}^{\rm PV}\right)^2
\le
\left(S_p^{\rm rated}\right)^2.
\]
Therefore
\[
Q_{p,t}^{\max}
=
\sqrt{
\left(S_p^{\rm rated}\right)^2-
\left(P_{p,t}^{\rm PV,av}\right)^2
}.
\]
At rated active power,
\[
Q_D^{\max}=\sqrt{231^2-210^2}\approx 96.23~{\rm kvar},
\]
and
\[
Q_H^{\max}=\sqrt{11^2-10^2}\approx 4.58~{\rm kvar}.
\]

\subsection{OpenDSS mapping}
The three-phase PV systems use line-to-line nominal voltage
\[
V_D^{\rm rated}=4.16~{\rm kV},
\]
whereas the single-phase wye systems use phase-to-neutral voltage
\[
V_H^{\rm rated}=2.4~{\rm kV}.
\]
No native \texttt{InvControl} object is introduced. The initial static operating
point uses zero reactive-power command. Later RBC, MPC, and RL controllers act
through the same OpenDSSSim requested-action interface.

\subsection{Fair-comparison consequence}
All P1 controller paradigms act on the same physical resources:
\[
\mathcal A^{P1}
=
\mathcal A_R^{\rm native}
\times
\mathcal A_C^{\rm native}
\times
\mathcal A_P^{10PV}.
\]
The feeder, DER placement, inverter ratings, native regulator/capacitor
equipment, disturbances, and physical bounds are therefore common across RBC,
MPC, and RL.
